\chapter{Conditional and Counterfactual Reasoning in Cognitive Tasks}
\label{sec:p19-rationality-6-3}

Byrne and Espino organize their chapter around four questions: whether
understanding a conditional involves possibilities or probabilities;
when people infer a conditional relation; whether factual and counterfactual
conditionals yield similar inferences; and how counterfactual negation is
represented. Their answer favors models of possibilities that can contain
symbolic annotations. That is an empirically motivated theoretical
position, not a theorem that every person uses one algorithm
\cite[chapter~6.3, pp.~405--418]{ByrneEspino2021ConditionalCounterfactual}.

The opening examples distinguish several tasks. A child given a
good-painting/treat conditional can expect a treat from the stated
antecedent. An office worker who misses the stipulated departure time
may predict missing dinner or search for another route. A learner driver
reflecting on an unsuccessful test may reason from a counterfactual
about having done everything right. These examples motivate inference
and imagination; they are not sampled observations with population
denominators. The chapter's later reaction-time and conclusion-production
experiments supply the actual evidence considered below.

\section{Inference patterns and competing representation accounts}
\label{sec:p19-rat63-accounts}

\begin{hbconvention}{def:p19-rat63-standing}{Language, formal comparisons and observations}
For the bounded formal comparisons, \(A,B\) are propositions, also
identified with events in one specified probability space when a
probability is written. The source's ``not-A'' and ``not-B'' are
written \(\neg A,\neg B\). The connective \(\to\) denotes
material implication only when that interpretation is explicitly stated.
An ordinary English conditional or counterfactual is not assigned that
meaning by its grammatical form alone. A schematic list of cognitive
models is a proposed representation, not measured brain content.
An ellipsis or unspecified component is not a false proposition.
\end{hbconvention}

The essay/mark example yields four traditional patterns:
\[
\begin{array}{c|c|c}
 \text{pattern}&\text{minor premise}&\text{proposed conclusion}\\\hline
 \text{modus ponens (MP)}&A&B\\
 \text{affirm consequent (AC)}&B&A\\
 \text{modus tollens (MT)}&\neg B&\neg A\\
 \text{deny antecedent (DA)}&\neg A&\neg B
\end{array}
\]
The source gives broad literature summaries: MP is usually highly
endorsed, whereas MT, AC and DA vary. Its approximate fractions are
not four counts from one experiment, nor universal rates. With fixed
material meanings MP and MT are valid and AC and DA are not, by
\cref{thm:p19-conditional-benchmark}. A cognitive account must explain
the observed pattern without treating an endorsement as proof of validity.

The chapter first describes abstract mental-rule accounts, where MP
has a direct rule but MT may need a derivation, and domain-sensitive
accounts using permission, obligation or social cost/benefit schemas.
The latter were particularly developed for selection tasks. It then
contrasts probabilistic/suppositional and mental-model approaches.
Suppositional accounts temporarily assume the antecedent and assess
the consequent, often using a conditional probability. Mental-model
accounts construct possibilities consistent with the interpreted assertion,
initially keeping few explicit cases and elaborating them when necessary.
Their models are proposed small-scale iconic simulations, not merely
rows of a written truth table. Parsimony limits the initially explicit
cases under working-memory constraints, and the principle of truth favors
cases consistent with the assertion rather than cases it rules out.
The account makes MP available from the first affirmative model, whereas
MT requires access to a negative-consequent case. The rule and probability
accounts also explain an MP/MT difficulty contrast in their own ways;
that shared prediction alone does not distinguish them. The model account's
no-counterexample standard concerns the interpreted possibilities, whose
selection can be modulated by meaning and pragmatic context. It is not
a claim that every English if-sentence has one fixed truth function.
These are families, not four completely specified algorithms.
\cite[Sec.~2, pp.~405--407]{ByrneEspino2021ConditionalCounterfactual}

For the apples/oranges example the model account proposes an initial
\(A\land B\) possibility and a marker that further cases may exist.
Some interpretations expand to a biconditional; a conditional
interpretation permits a third case. The two proposed fully elaborated
comparisons are:
\[
\begin{array}{c|l|l}
 &\text{fully elaborated possibility reading}
 &\text{antecedent-focused proposed comparisons}\\\hline
 \text{conditional}&AB,\ \neg A\neg B,\ \neg AB
                   &AB,\ A\neg B\\
 \text{biconditional}&AB,\ \neg A\neg B
                     &AB,\ A\neg B,\ \neg AB
\end{array}
\]
Juxtaposition means conjunction. The second column is not automatically
the person's initial representation; the third is an algorithmic hypothesis
about the cases considered, not an axiom of probability theory.
For example, \(P(\neg B\mid A)=1-P(B\mid A)\) does not logically
require an explicit mental image of an \(A\neg B\) case.
Predictions about reading times require an additional account of access
to representations and of how access affects the recorded response.

\section{Possibility priming and inferring a conditional relation}
\label{sec:p19-rat63-factual-evidence}

\subsection{A conditional versus a biconditional as a prime}

In the source's market narrative, Carmen reads a sign about apples and
oranges and then encounters a statement about the shelves. The critical
target is either an affirmative conjunction or a negated-antecedent
conjunction. The model account predicts that \(\neg A\land B\)
is more accessible after a conditional than after a biconditional because
it is compatible with the first reading but not the second. The
antecedent-focused account considered by the authors instead treats a
biconditional as two directed conditionals, making that case available
through its converse. These predictions depend on that particular
representation and latency mapping; they are not derived from Bayesian
coherence alone.
\cite[Sec.~2.1, pp.~407--409]{ByrneEspino2021ConditionalCounterfactual}

\paragraph{The selected original task and sample.}
Espino, Santamar\'{i}a and Byrne's Experiment~2 used a new cohort of
32 University of La Laguna students for course credit. They read
seven-sentence narratives in their native Spanish, one sentence at a
time, advancing by keypress. The fully within-person design crossed
conditional/biconditional primes with affirmative/negated-antecedent
targets: 12 instances per cell, 48 experimental narratives, plus
32 fillers, in a different randomized order for each participant.
Simple comprehension questions required yes/no responses and did not
ask participants to produce a logical conclusion. Six practice
scenarios preceded testing. Age/sex distributions and recruitment
dates are not specified in the selected report.
\cite[pp.~1074--1077]{EspinoSantamariaByrne2009True}

\paragraph{Measurement, preprocessing and result.}
The measure was target reading latency. The inherited rule replaced
latencies below half the mean or above the mean plus 2.5 SD with
those respective boundary values; 3\% were replaced in Experiment~2.
Only participants with more than 85\% comprehension accuracy and
latencies associated with correct responses entered the reported
analysis. The report does not give a separate exclusion count to
add to its 32-participant sample.
The negated-antecedent means were 1,894~ms (SD~382) after a
conditional and 2,157~ms (SD~701) after a biconditional. Their rounded
table difference is 263~ms; the prose reports 262~ms. The planned
comparison was \(F(1,31)=8.54,p<.006\), and the interaction
\(F(1,31)=5.62,p<.025\). Affirmative targets were 1,451 and
1,487~ms, with no reliable 36-ms difference. These are condition
means and SDs, not a confidence interval or independent observations
equal to all repeated trials.

The result supports the tested accessibility contrast. It does not
establish that no probabilistic representation was used, exclude every
alternative linguistic-processing explanation, or identify a neural
format. The chapter interprets it more broadly as evidence for thinking
about true possibilities. That interpretation must retain its tested
conditions and competing accounts' actual algorithmic commitments.

\subsection{Garden disjunctions and the compatibility heuristic}

The source next asks whether roses-or-lilies supports a conditional
from no roses to lilies, and reverses the direction using a
no-tulips/daffodils conditional. Negated clauses provide further cases,
such as no-daisies-or-no-buttercups and a daisies/no-buttercups
conditional. This is not necessarily a causal relation.
For material interpretation, the corresponding formulas have identical
truth possibilities. For a conditional-probability assessment, the
one-premise confidence comparison is different.

\begin{hblemma}{thm:p19-rat63-or-if}{proved=annex:p19.rat63-or-if}{Truth equivalence and the source's probability comparison}
Material \(A\lor B\) and \((\neg A)\to B\) have the same
truth set, so entail each other. On a law with \(P(\neg A)>0\),
put \(q=P(B\mid\neg A)\). Then
\[
 P(A\lor B)=P(A)+(1-P(A))q,
 \qquad P(A\lor B)-q=P(A)(1-q)\geq0.
\]
If one-premise p-validity requires conclusion confidence to be at
least premise confidence under every such law, the inference from
degree \(q\) to disjunction probability is p-valid, but the reverse
need not be. A four-cell law \((.4,.4,.04,.16)\), in order
\((AB,A\neg B,\neg AB,\neg A\neg B)\), gives disjunction
probability \(.84\) and conditional degree \(.2\).
If the disjunction is certain and \(P(\neg A)>0\), however,
\(q=1\). The failure of confidence nondecrease is not a
counterexample to that certain-premise conclusion.
\end{hblemma}

Thus the same word ``valid'' cannot silently switch between truth
preservation, confidence nondecrease and an experimental endorsement.
The probability degree is not automatically the probability of the
material-implication event. The related Equation and fixed-event
distinction are \cref{def:p19-rat46-equation}.

The model account also proposes a shortcut: inspect an available
compatible possibility without exhaustively searching for an incompatible
premise case. Espino and Byrne call this a compatibility heuristic.
It can accept a conclusion that is possible rather than necessary.
Initial and elaborated representations must not be conflated; an
unspecified component of an initial model is not an explicit negation.
The procedure first constructs the initial conditional and disjunctive
models, then looks for a compatible case. If that fails, it can construct
one immediately available alternative for the conditional and try again.
The proposed control assumptions limit this search to one alternative,
alter the conditional rather than elaborating the disjunction, and
obtain the alternative by negating the conditional's two mentioned
components. Thus \(A\to B\) and \(\neg A\lor\neg B\) can receive
an erroneous endorsement after the shared alternative \(\neg A\neg B\)
is considered. Conversely, the valid material inference from
\(A\to B\) to \(\neg A\lor B\) can be missed when incomplete
initial representations do not make a shared case explicit. Further
deliberation can repair that response; the heuristic is not an assertion
that every participant stops at those steps.
\cite[Secs.~2.1--2.3, pp.~100--108]{EspinoByrne2013Compatibility}

\begin{hblemma}{thm:p19-rat63-compatibility}{proved=annex:p19.rat63-compatibility}{Compatibility does not establish entailment}
Let \(\mathcal T(\varphi)\) be a formula's set of two-proposition
truth assignments. For \(P_0=A\lor B\) and material
\(Q_0=A\to B\), \(\mathcal T(P_0)\cap\mathcal T(Q_0)\)
is nonempty, but \(\mathcal T(P_0)\nsubseteq\mathcal T(Q_0)\).
The assignment \(A=1,B=0\) is a counterexample to necessity.
Finding a mutually compatible case therefore cannot substitute for
checking that all premise cases satisfy the conclusion.
\end{hblemma}

\paragraph{The selected original inference task.}
Espino and Byrne's Experiment~1b used 60 La Laguna undergraduates,
none reporting formal logic training, tested in groups with self-paced
booklets. It supplied 32 problems: four conditional polarities crossed
with four inclusive-disjunction polarities, in both directions, using
animals, fruit, flowers and kitchen-implement contents. Participants
were instructed to judge necessity given the premise and could answer
yes, no, or not possible to know. The two direction blocks were
counterbalanced and their internal order randomized. This differs
from Experiment~1a's separate cohort and binary response format.
\cite[Secs.~3.1--3.2, pp.~109--112]{EspinoByrne2013Compatibility}

Compatibility categories differed: endorsement was 62\% for a
compatible initial possibility, 46\% for a compatible alternative,
26\% for shared first polarity only, and 30\% for shared second
polarity only. The corrected compatibility ANOVA was
\(F(2.204,127.854)=27.49,p<.001\).
For the affirmative-disjunction transformation highlighted by the
2021 chapter, ``or therefore if'' was endorsed by 48\% versus
28\% for its corresponding reverse. Across all 32 forms, however,
the direction averages were 40\% and 42\%, with no significant
main effect, \(F(1,58)=1.34,p=.25\). The result for the highlighted
form must not be generalized to every conditional/disjunctive inference.
No confidence interval for its 20-percentage-point contrast is supplied
in the selected prose; the figures' error bars are SEM. Raw individual
responses were not reanalyzed here.

The source authors regard the chosen asymmetry as challenging the
probability account's predictions and as supporting possibility-based
reasoning. The original paper also emphasizes widespread compatibility
errors rather than systematic full truth-table deduction. Neutral
content reduces some prior-belief explanations, but response instructions,
negation processing and the exact procedure linking confidence or
compatibility to a yes judgment remain part of the comparison.

\section{Counterfactual conjectures and the facts they presuppose}
\label{sec:p19-rat63-counterfactuals}

The source's locational example contrasts an indicative conditional
linking Paolo's presence in Venice with Marco's presence in Padua and
its subjunctive counterpart. Given Marco's absence, the reviewed studies
report MT more readily for the counterfactual. DA is also reported more
often, whereas MP and AC are reported similarly for the two moods.
Figure~6.3.3 extends the reviewed MT contrast across locational, causal,
definitional, everyday, promise, threat and deontic contents. These
bars draw on different studies, not one factorial experiment or a pooled
sample with one denominator. The chapter's description of a robust
effect retains a review-level boundary here; recruitment and precision
cannot be supplied by inventing records for those unselected studies.
\cite[Sec.~3.1, pp.~410--413]{ByrneEspino2021ConditionalCounterfactual}

The model account explains the contrast through the initially available
possibilities. For an ordinary counterfactual about apples and oranges,
its proposed initial representation has two labeled rows:
\[
\begin{array}{c|cc}
 \text{counterfactual conjecture}&A&B\\
 \text{presupposed facts}&\neg A&\neg B
\end{array}
\qquad\ldots
\]
The labels matter. These are not two assertions conjoined as facts in
one world. The first represents an alternative to the second. Under a
fully elaborated conditional interpretation, a further \(\neg AB\)
counterfactual possibility can be included. The theory therefore does
not identify every counterfactual with a biconditional or deny contextual
variation. The fully elaborated state sets can parallel those of factual
conditionals while carrying different epistemic-status labels.

On this account, a \(\neg B\) minor premise matches an already
available negative row, so MT does not first require elaborating that
row as it would for a parsimonious factual representation. An affirmative
\(A\) premise can instead cause the proposed process to discard the
negative row and retag the affirmative row as factual, supporting MP.
This is a model of how a participant integrates task information, not a
license to conjoin mutually inconsistent factual premises. Nor does the
ease of DA make DA truth-preserving under the fixed material semantics
of \cref{thm:p19-conditional-benchmark}.

The source adduces memory intrusions, inferred implications, priming,
eye movements and ERP/fMRI measures in support of ready access to both
conjecture and presupposed facts. These are different operational
measures. Their chapter-level convergence claim is not an independently
reconstructed neural mechanism here. The selected original priming
experiment permits a more exact account of one part of the evidence.
Specifically, the reviewed memory task found mistaken recollection of
the negative facts as if they had been explicitly stated; an implication
judgment asked what a speaker meant to convey. The visual-world account
reports more looks at the negative-fact images after a counterfactual
than after a factual conditional. These observations constrain access
and interpretation hypotheses, not the claim that a literal negative
sentence was spoken. The cited ERP and fMRI studies use physiological
measures with different tasks; their mention does not independently
identify the contents or format of every representation.
\cite[p.~412]{ByrneEspino2021ConditionalCounterfactual}

\subsection{The selected factual/counterfactual priming experiment}

Santamar\'{i}a, Espino and Byrne's Experiment~1 tested 32 La Laguna
students for course credit. Participants read native-Spanish narratives
sentence by sentence and advanced by keypress. The within-person design
crossed factual versus counterfactual primes with all four target
conjunctions: \(AB\), \(\neg A\neg B\), \(\neg AB\), and
\(A\neg B\). Six instances per cell gave 48 experimental trials;
32 fillers and six practice narratives accompanied them. Contents were
randomly assigned and trials randomized for each participant. The
indicative primes used present tense and the counterfactual primes used
subjunctive past tense. The comparison consequently changes a linguistic
package, not an isolated neural representation while holding every other
feature fixed. Comprehension questions did not request an inference.
\cite[pp.~1150--1151]{SantamariaEspinoByrne2005Alternatives}

The recorded measure was target reading latency. The same boundary
replacement rule used in the later 2009 study replaced 5\% of latencies
in this experiment. Only participants exceeding 90\% comprehension
accuracy and latencies associated with correct answers entered analysis;
a separate recruited-versus-retained exclusion count is not supplied.
Table~1 reports the following means in milliseconds, with SDs in
parentheses:
\[
\begin{array}{c|rrrr}
 \text{prime}&AB&\neg A\neg B&\neg AB&A\neg B\\\hline
 \text{factual}&1646\ (272)&1838\ (314)&2182\ (478)&1928\ (458)\\
 \text{counterfactual}&1702\ (335)&1724\ (308)&2244\ (490)&2055\ (421)
\end{array}
\]
The planned negative-conjunction contrast was 114~ms faster after
the counterfactual, \(t(31)=2.34,p<.03\), reported standard error
48.66~ms. The affirmative contrast was not significant,
\(t(31)=1.26,p=.21\); neither was the \(\neg AB\) contrast
(\(p=.50\)) or the \(A\neg B\) contrast (\(p=.09\)).
The omnibus conditional-by-conjunction interaction had \(p=.079\),
not significance at .05, and the conditional main effect had \(p=.34\).
The authors base their planned claim on the predicted negative contrast;
it must not be retold as a significant omnibus interaction. Absence of
a significant affirmative contrast is not a demonstrated equality of
processing times. No participant-level interval or reanalysis is added.

The task supports a selective accessibility difference under these
primes. It does not establish that every counterfactual explicitly
asserts \(\neg A\land\neg B\), or that faster reading uniquely
identifies an iconic format. The selected study's semifactual experiments
are outside this numerical record; they are not silently counted as
additional replications of this particular factual/counterfactual contrast.

\subsection{Competing accounts and the temporal comparison}

The source considers a suppositional account that combines a prior
conditional assessment with beliefs about implied facts into one degree
of belief, while favoring a single considered situation. It asks how that
account supplies ready access to both labeled possibilities and the
inference difference. This is an objection to specified algorithmic
commitments, not a deduction that any probability calculation must
exclude representations of facts or alternatives.

The Kennedy examples make the background facts and temporal question
consequential. An indicative inference after a shooting has already
occurred asks who actually fired; a counterfactual about removing
Oswald as shooter can ask whether anyone would have fired at all.
The source's suggested parallels align a future-oriented indicative
question with the uncertain counterfactual, and align a retrospective
counterfactual identifying an already-presupposed shooter with the
retrospective indicative. The original pair's contrasting plausibility
does not isolate grammatical mood while preserving these presuppositions
\cite[p.~413]{ByrneEspino2021ConditionalCounterfactual}.

The causal-Bayes-net account introduces a different contrast: observing
a cause absent versus preventing it. The chapter reviews moving-object
tasks where the proposed MT response differs between observational and
preventive questions. Those questions need not have the same formal
target; changing a causal equation is not merely adding an observation.
The chapter also objects that causal prevention alone does not explain
the broader noncausal and deontic inference pattern. Neither objection
is answered by replacing all cognitive tasks with a structural causal
model.

\begin{hbconvention}{def:p19-rat63-world-typing}{Keep the factual evidence and queried world distinct}
In the P16 semantics of \cref{def:p16-counterfactual-worlds},
\(C[B=0]\) denotes the value of \(C\) after replacing \(B\)'s
equation, for the same background unit. If factual evidence is \(B=1\)
and has positive probability, a typed counterfactual query is
\(P(C[B=0]=c\mid B=1)\). It is not conditioning on the impossible
same-world conjunction \(B=1\land B=0\). The abduction/action/
prediction operation in \cref{def:p16-counterfactual-evidence} keeps
the factual evidence fixed while evaluating the changed equations.
\end{hbconvention}

\begin{hblemma}{thm:p19-rat63-observe-intervene}{proved=annex:p19.rat63-observe-intervene}{A worked distinction, not a cognitive model}
Let \(U\in\{0,1\}\) have probability \(1/2\) at each value,
and let the factual equations be \(B=U\), \(C=B\).
Then \(P(C=0\mid B=1)=0\), but
\(P(C[B=0]=0\mid B=1)=1\). Ordinary conditioning on
\(B=1\land B=0\) is undefined under the positive-evidence convention.
\end{hblemma}

This illustration justifies the typing distinction in the chapter's
mixed factual/intervention notation; it does not predict the endorsement
rate of a participant in a prevention task. Cognitive representation,
causal truth in a stipulated model and identifying a causal model from
observations remain separate questions.

The false-antecedent and false-consequent implications used in the dual
row example are contextual presuppositions. An Anderson conditional about
arsenic and the observed symptoms can use subjunctive wording while
leaving it open that arsenic was actually taken. Mood alone therefore
does not establish counter-to-fact status. For a genuinely impossible
antecedent, such as the source's squaring-the-circle counterpossible,
which asks whether Hobbes achieving that impossibility would make ill
children in South America care, the psychological evidence discussed
is scant; the ordinary examples do not settle its semantics.
Material implication's automatic truth
at a false antecedent is precisely not an adequate test of every
counterfactual's informativeness.

\section{Negation as an alternate or a symbolic annotation}
\label{sec:p19-rat63-negation}

For the flowers/trees counterfactual, the conjecture is roses with
orange trees. Its presupposed opposite can be represented as no roses
with no orange trees, or, in a suitable binary context, as poppies with
apple trees. The embodied specific-alternates proposal connects meaning
to modality-specific sensory or motor simulations and represents negation
through an alternate, sometimes after first representing and inhibiting
the affirmative. The annotated-model account permits iconic simulations
to contain a symbolic negation tag. Its claim is plural: context and
resource demands can change the representation
\cite[Sec.~3.2, pp.~413--415]{ByrneEspino2021ConditionalCounterfactual}.

In a binary flower context, both proposals predict a conclusion naming
poppies from a negative minor premise about orange trees, rather than
a conclusion saying no roses. In a multiple context, adding lilies
makes a negative flower category compatible with more than one named
alternate. The specific-alternates view predicts continued enumeration,
whereas the annotated view allows the compact negative representation.
The source treats the switch as a response to working-memory demands;
the selected experiments manipulate context, not working-memory capacity
directly. That proposed mechanism must not be turned into an independently
measured mediator.

\begin{hblemma}{thm:p19-rat63-alternates}{proved=annex:p19.rat63-alternates}{What a supplied exhaustive category context permits}
For a single flower category whose mutually exclusive and exhaustive
options are roses or poppies, each possible, \(\neg\text{roses}\)
is equivalent to poppies. If the options are roses, poppies or lilies
instead, each possible,
\(\neg\text{roses}\) is equivalent to poppies-or-lilies but does
not determine either alternate individually. Without exhaustiveness,
the negation need not imply even the disjunction of the named alternates.
These equivalences do not license DA or AC from a material conditional.
\end{hblemma}

Thus affirmative and negative wording can encode the same category
information under the stipulated context. A preference for one wording
is not necessarily a difference in logical accuracy. The experiment's
claim that wording reflects the representation used to derive it is an
additional measurement assumption, acknowledged by its authors. Output
wording could also depend on response selection or linguistic habits.

\subsection{The selected binary- and multiple-context experiments}

Espino and Byrne's Experiment~1 tested 32 psychology undergraduates at
La Laguna, without formal logic training or previous reasoning-experiment
participation. The source describes the recruiting student pool generally
as aged 18--24 and about two-thirds women; this is not an exact sex/age
tabulation of the recruited 32. In small groups participants completed
self-paced booklets, generating a necessary conclusion or stating that
there was none. A within-person \(4\times2\) design crossed MP, DA,
AC and MT with direct versus indirect reference; two instances per cell
gave sixteen problems with randomly assigned contents and counterbalanced
term order. Participants were told not to revisit completed answers.
The English examples reproduce the paper's illustrative and translated
presentation; they are not a claim that those exact English words were
the administered stimulus transcripts.
\cite[EarlyView pp.~11--15]{EspinoByrne2018Negation}

Direct reference uses the terms in the counterfactual itself: roses for
MP, and no orange trees for MT. Indirect reference uses their contextual
alternates: no poppies for MP, and apple trees for MT. It therefore
separates the inference type from whether the minor premise happens
to be verbally negative. Conclusions were coded as affirmative or
negative standard endorsements under those category alternatives.
The study did not record motor compatibility or brain activation in
these experiments.

Experiment~1 produced more affirmative than negative conclusions in
each of the eight inference/reference cells, with the reported binomial
sign tests \(p<.0001\). For negative-premise direct MT and DA, the
affirmative percentages were 84\% and 86\%; for indirect reference
they were 91\% and 97\%. The four direct-versus-indirect comparisons
were not significant at the corrected \(\alpha=.0125\). The result
supports an alternates tendency in this binary task; it does not yet
distinguish the two representation accounts, which both predicted it.

Experiment~2 recruited a different 32 La Laguna undergraduates without
formal logic training and used the same method. It added a third
category on the side queried by the conclusion: a multiple antecedent
for AC/MT, or a multiple consequent for MP/DA. It did not make both
sides multiple simultaneously. In multiple contexts, a complete
affirmative endorsement named the disjunction of the alternatives;
the coding also recorded a partial affirmative response naming just
one alternate. These are the administered and coded forms, not a
universal experiment on every possible negation
\cite[EarlyView pp.~16--21]{EspinoByrne2018Negation}.

The result is more qualified than ``no alternates in multiple contexts.''
For direct DA and MT, affirmative responses were 50\% and 48\%; the
affirmative-versus-negative sign tests were \(p=.61\) and \(p=1.0\).
Indirect DA and MT retained more affirmative than negative responses,
both \(p<.0001\), with 66\% affirmative responses for each. Direct
MP/AC also retained the affirmative tendency. The inference-by-reference
interaction was \(F(1.599,49.583)=10.11,p<.002\), with the
Greenhouse--Geisser correction and reported \(\eta_p^2=.25\).
Although the direct/indirect DA and MT comparisons had uncorrected
\(p<.025\) and \(p<.02\), neither meets the stated Bonferroni
threshold .0125. They must not be reported as significant corrected
contrasts. The primary explicitly calls the direct negative-premise
effect reduced, not reversed, and notes heterogeneity across participants
and trials rather than a uniform strategy switch.

The article's public-data references do not constitute a participant-level
reanalysis here. No new interval, pooled effect across its nine experiments,
or proof of a sensory versus symbolic neural format is supplied. The two
selected experiments support a context-sensitive output pattern compatible
with the annotated account; the inference from that pattern to an internal
format retains the stated measurement assumption and alternative response
mechanisms. Subsequent experiments in the article are not automatically
added to this fixed numerical record.

\section{What the four questions establish}

Byrne and Espino conclude that possibility-based representations connect
conditional inference with imagination, that counterfactuals make both
conjectures and presupposed facts accessible, and that models can mix
iconic and symbolic elements. The four selected evidence records show
why these are substantive proposals: they predict differences in priming,
inference direction or conclusion wording under specified tasks. They
also show why the conclusion must remain bounded. Latency and endorsement
are not direct readouts of a unique algorithm; an aggregate response
pattern is not a normative instruction; and a probability law does not
itself determine an algorithmic representation or a response frequency.

The related \cref{sec:p19-rationality-4-6} develops Oberauer and
Pessach's differently argued probability interpretation, including
supposition and the Equation. The present source's possible-versus-
probable conclusion and that probabilistic argument should remain
distinct, not be declared reconciled by changing the meaning of
``conditional.'' \Cref{sec:p19-conditional-reasoning-suppression,
sec:p19-conditional-interpretation-experiments} treats extra premises
and changing interpretation, not the same priming or alternate-output
task. The simulation/emotion material in
\cref{sec:p19-rationality-2-4} concerns another operational question;
an Olympic rating is not a conditional-inference measure. P16 supplies
formal cross-world queries, not an empirical theory of sentence
comprehension. The evidence distinctions of
\cref{sec:p19-evidence-methods} govern all of these comparisons.

The practical consultation rule is therefore to identify the question
first: truth-preserving inference under a declared interpretation,
confidence comparison under a probability assessment, comprehension
priming, generated conclusion, or a causal intervention query. A success
under one does not settle the others. The source's contribution is its
explicit set of contrasts and experimental predictions, with their
representation hypotheses and objections kept visible, not a replacement
of conditional reasoning by a single truth table, probability number
or structural equation.
